%% shortname: Aggregate Demand, Wealth Distribution and WARP
\documentclass[11pt]{article}
\usepackage{microstyle}

\begin{document}

\lessonheader{11}{Chapter 4: aggregation across consumers --- aggregate demand, aggregate wealth and WARP}{}


\section{A small detour: aggregation across consumers}

A small detour to study aggregation across consumers. It matters for general
equilibrium, macroeconomics, and empirical work.

\runin{Main objective.} Be cautious about reading too much into aggregate
behavior, especially in the context of welfare.

We consider three explicit questions.
\begin{flatlist}
  \item The first and second questions are about whether aggregate demand
        ``looks like'', or has the same properties as, individual demand.
  \item The third question asks whether aggregate demand is useful for welfare
        evaluation.
\end{flatlist}
In all three cases, strong restrictions are needed to arrive at the answers we
hope for.

\section{Aggregate demand and aggregate wealth}

When does aggregate demand look like individual demand, in the sense that it
depends only on prices and aggregate wealth? In other words, does the
distribution of wealth matter too?

The answer to this first question is generally \emph{no}, unless more
restrictions are imposed.

\runin{MWG notation.} Let $x_i(p,w_i)$ denote agent $i$'s Marshallian (Walrasian)
demand, and $w_i$ her wealth. So $x_i(p,w_i)\in\R^L_+$ is a vector, i.e.\ a
bundle. Agent $i$'s demand for good $k$ is $x_{ki}(p,w_i)$.

With $I$ agents, aggregate demand is
\[ x(p,w_1,\dots,w_I)=\sum_{i=1}^{I}x_i(p,w_i). \]

When does
\[ x(p,w_1,\dots,w_I)=x\Bigl(p,\sum_{i=1}^{I}w_i\Bigr)\,? \]
This is not easy, because there is a loss of information from ignoring the
distribution.

\topic{What does it take for an affirmative answer?}
If it holds, then $x\bigl(p,\sum_i w_i\bigr)$ should be the same even after a
redistribution of wealth that leaves $\sum_i w_i$ unchanged. So examine wealth
changes $(dw_1,\dots,dw_I)$ with $\sum_{i=1}^{I}dw_i=0$. Taking the total
derivative, for every good $k$,
\[ \sum_{i=1}^{I}\pd{x_{ki}(p,w_i)}{w_i}\,dw_i=0. \]
This holds for all $(w_1,\dots,w_I)$ and all such wealth changes if and only if,
when comparing any two agents $i$ and $j$,
\[ \pd{x_{ki}(p,w_i)}{w_i}=\pd{x_{kj}(p,w_j)}{w_j}
   \qquad\text{for every good }k. \tag{$*$} \]
(For ``only if'', take $dw_i=1$, $dw_j=-1$ and all other changes zero.)

\begin{flatlist}
  \item That is, wealth effects are the same for all agents.
  \item Moreover, wealth effects must be constant across wealth levels,
        i.e.\ independent of $w_i$ (($*$) must hold whatever $w_i$ and $w_j$ are).
  \item In fact, individual demand must be linear (affine) in $w_i$: the wealth
        expansion paths are parallel straight lines.
\end{flatlist}
Sounds restrictive.

\topic{Two cases where ($*$) holds}
Condition ($*$) is satisfied in two cases.
\begin{flatnum}
  \item All agents have preferences that are quasilinear with respect to the same
        good. With $L=2$ and good 2 as the num\'eraire,
        $u_i(x)=x_2+\phi_i(x_1)$: the demand for good 1 does not depend on
        wealth, and all extra wealth is spent on good 2 (wealth effects $0$ and
        $1/p_2$, the same for everyone).
  \item All agents have identical preferences that are homothetic: there is a
        common utility function that is homogeneous of degree 1,
        \[ u(\alpha x)=\alpha\,u(x),\qquad \alpha>0. \]
        The trade-off between two goods depends only on the ratio $x_1/x_2$, so
        demand is $x_i(p,w_i)=w_i\,x(p,1)$ and wealth effects are constant.
\end{flatnum}

\begin{center}
\begin{tikzpicture}[scale=1.15]
  % homothetic: u = sqrt(x1 x2), p = (1,1), w = 2 and w = 3
  \ecaxes{3.4}{3.4}
  \draw[budget] (0,2) -- (2,0);
  \draw[budget] (0,3) -- (3,0);
  \draw[domain=0.36:2.8,smooth,samples=60] plot (\x,{1/\x});
  \draw[domain=0.78:3.2,smooth,samples=60] plot (\x,{2.25/\x});
  \draw[dashed] (0,0) -- (3.0,3.0) node[ecnode,above right]{$x_1/x_2$ constant};
  \node[bundle] at (1,1) {};
  \node[bundle] at (1.5,1.5) {};
  \node[ecnode] at (1.7,3.75) {Homothetic};
  \node[ecnode,below] at (3.4,0) {$x_1$};
  \node[ecnode,left]  at (0,3.4) {$x_2$};
  \node[ecnode,align=left,anchor=west] at (3.0,1.3) {wealth effects\\ are constant};
\end{tikzpicture}
\hspace{1.2cm}
\begin{tikzpicture}[scale=0.85]
  % quasilinear: u = x2 + 2 ln x1, p = (1,1), w = 3 and w = 4
  \ecaxes{4.6}{4.6}
  \draw[budget] (0,3) -- (3,0);
  \draw[budget] (0,4) -- (4,0);
  \draw[domain=0.75:3.4,smooth,samples=60] plot (\x,{2.386-2*ln(\x)});
  \draw[domain=0.75:4.3,smooth,samples=60] plot (\x,{3.386-2*ln(\x)});
  \draw[dashed] (2,0) -- (2,4.2);
  \node[bundle] at (2,1) {};
  \node[bundle] at (2,2) {};
  \node[ecnode] at (2.3,5.0) {Quasilinear: $u=x_2+\phi(x_1)$};
  \node[ecnode,below] at (4.6,0) {$x_1$};
  \node[ecnode,left]  at (0,4.6) {$x_2$};
\end{tikzpicture}
\end{center}

\topic{The Gorman form}
In both cases, the indirect utility has the \emph{Gorman form}
\[ v_i(p,w_i)=a_i(p)+b_i(p)\,w_i. \]
In fact, ($*$) holds if and only if every agent has an indirect utility of the
Gorman form with
\[ b_i(p)=b(p),\qquad\text{the same for all agents.} \]
(Homothetic case: $v(p,w_i)=w_i\,v(p,1)$, so $a_i=0$ and $b(p)=v(p,1)$.
Quasilinear case: $b(p)=1/p_2$.)

\topic{Conclusion so far}
It might work, but we need strong assumptions on how alike agents are.

But we asked a hard question: we wanted it to work for \emph{every} wealth
distribution. This ignores what we might know about the wealth distribution.
Maybe we should be glad if it works for a nice subset of wealth distributions.

\topic{A deterministic wealth distribution rule}
Let $w$ denote aggregate wealth, and imagine that there is a deterministic wealth
distribution, i.e.\ rules in place (e.g.\ through government redistribution, or
the agents' ownership of assets in this economy). Agent $i$'s wealth is
$w_i(p,w)$. Then aggregate demand is a function of prices and aggregate wealth:
\[ \sum_{i=1}^{I}x_i\bigl(p,w_i(p,w)\bigr)=x(p,w). \]

\section{Aggregate demand and WARP}

When does aggregate demand $\sum_i x_i(p,w_i)$ satisfy WARP?

Remember that WARP is a minimal consistency assumption on demand at the
individual level. Aggregate demand satisfies continuity, homogeneity of degree
zero and Walras' law, because individual demands have all these properties.
What about WARP?

For the best chance, assume there is a deterministic wealth distribution rule,
and that it takes the form
\[ w_i(p,w)=\alpha_i w,\qquad \alpha_i>0,\quad \sum_{i=1}^{I}\alpha_i=1. \]
Aggregate demand $x(p,w)$ satisfies WARP if, for any two price--wealth
situations $(p,w)$ and $(p',w')$,
\[ p\cdot x(p',w')\le w\ \text{ and }\ x(p',w')\ne x(p,w)
   \;\Longrightarrow\; p'\cdot x(p,w)>w'. \]

\topic{A counterexample}
Even in this scenario, there are counterexamples: WARP may fail in the
aggregate. Take $I=2$, $\alpha_1=\alpha_2=\tfrac12$, so both agents have the same
budget set, and $L=2$.

\begin{center}
\begin{tikzpicture}[scale=2.0]
  \ecaxes{3.9}{3.3}
  % B_{p,w/2}: x2 = 3 - (5/3) x1 ;  B_{p',w'/2}: x2 = 1.8 - 0.5 x1
  \draw[budget] (0,3) -- (1.8,0);
  \draw[budget] (0,1.8) -- (3.6,0);
  \node[ecnode,right] at (0.15,2.85) {$B_{p,\,w/2}$};
  \node[ecnode,above right] at (2.6,0.6) {$B_{p',\,w'/2}$};
  % agent 1 and agent 2 demands
  \node[bundle] at (0.6,2.0) {};   \node[ecnode,right] at (0.6,2.05) {$x_1$};
  \node[bundle] at (1.7,1/6) {};   \node[ecnode,right] at (1.7,0.2) {$x_2$};
  \node[bundle] at (0.1,1.75) {};  \node[ecnode,above right] at (0.08,1.75) {$x_1'$};
  \node[bundle] at (1.2,1.2) {};   \node[ecnode,above right] at (1.2,1.2) {$x_2'$};
  % averages (each lies on its own budget line)
  \node[bundle,fill=red!75!black] at (1.15,13/12) {};
  \node[ecnode,below left,text=red!75!black] at (1.15,1.08) {$A$};
  \node[bundle,fill=red!75!black] at (0.65,1.475) {};
  \node[ecnode,below left,text=red!75!black] at (0.67,1.49) {$B$};
\end{tikzpicture}
\end{center}

Here $x_i$ is agent $i$'s demand at $(p,w/2)$ and $x_i'$ her demand at
$(p',w'/2)$.
\begin{flatlist}
  \item Agent 1 prefers $(p,w)$: $x_1'$ was affordable at $(p,w/2)$, but $x_1$ is
        not affordable at $(p',w'/2)$.
  \item Agent 2 prefers $(p',w')$: $x_2$ was affordable at $(p',w'/2)$, but
        $x_2'$ is not affordable at $(p,w/2)$.
  \item Each agent's choices satisfy WARP.
  \item $A=\tfrac12(x_1+x_2)$ is average demand given $(p,w)$;
        $B=\tfrac12(x_1'+x_2')$ is average demand given $(p',w')$.
\end{flatlist}
$B$ is affordable at $(p,w)$ when $A$ is chosen, and $A$ is affordable at
$(p',w')$ when $B$ is chosen. It looks as if $A$ is revealed preferred to $B$,
but $B$ is also revealed preferred to $A$ $\Rightarrow$ aggregate demand violates
WARP.

\topic{Conclusion}
Disappointing: the answer to our question is no. Aggregate demand is not
guaranteed to satisfy WARP. The next best thing: under what conditions is the
answer affirmative instead?

On the other hand, even if aggregate demand violates WARP, individual agents can
still be rational.

\begin{transcriptnotes}
  \item This replaces the earlier one-page transcription; the full original is
        9 pages. Page 1 is a cover page reading ``Chapter 4'' only; the subtitle
        is written to match the earlier lessons. No date is given.
  \item Notation follows Lessons 7--10: the original's capitals $X$, $P$, $W_i$,
        $V$, $A_i$, $B_i$ are typeset as $x$, $p$, $w_i$, $v$, $a_i$, $b_i$.
        Agent $i$'s demand for good $k$ keeps the original's index $k$ (MWG
        writes $x_{\ell i}$). The label ($*$) for the condition is the
        original's; the original refers to ``Condition ($*$)'' without marking
        the equation, so it is attached to the equal-wealth-effects condition.
  \item \textbf{Corrected:} the Gorman form is written $V_i(P,u_i)$ in the
        original; indirect utility is a function of prices and wealth, so it is
        $v_i(p,w_i)$.
  \item \textbf{Corrected:} ``There's a common wealth function that are
        homogeneous degree 1, $V(\alpha x)=\alpha V(x)$'' --- it is a common
        \emph{utility} function of bundles $x$; written $u(\alpha x)=\alpha
        u(x)$.
  \item \textbf{Corrected:} the quasilinear utility is written ``$x_2+F(x_2)$'';
        the function must be of the other good, $x_2+\phi(x_1)$. The figure (a
        vertical dashed line through the tangencies) shows demand for good 1
        fixed, so good 2 is the num\'eraire.
  \item \textbf{Added:} the reason for ``only if'' in ($*$); the line ``wealth
        expansion paths are parallel straight lines''; the demand
        $x_i=w_i\,x(p,1)$ in the homothetic case; the wealth effects $0$, $1/p_2$
        in the quasilinear case; the Gorman-form parameters for both cases; the
        statement of WARP for aggregate demand (MWG Definition 4.C.1); and the
        explanation of why each agent satisfies WARP and why $A$, $B$ violate it.
        In the quasilinear case the wealth effects are $0$ and $1/p_2$ only for
        interior solutions (wealth large enough), as in MWG.
  \item Figures are redrawn with specific numbers so the geometry is exact:
        homothetic $u=\sqrt{x_1x_2}$; quasilinear $u=x_2+2\ln x_1$; in the WARP
        figure, $B_{p,w/2}\colon x_2=3-\tfrac53x_1$ and
        $B_{p',w'/2}\colon x_2=1.8-0.5x_1$. In the WARP figure the agent labels
        $x_1,x_2,x_1',x_2'$ are the original's (subscripts are agents, not
        goods). The original's ``1 prefers $(p,w)$, 2 prefers $(p',w')$'' and the
        red notes are set as the bullet list below the figure.
  \item Page 2: ``aggregation cross consumers'' $\to$ \emph{across consumers};
        ``be cautious of reading to much into'' $\to$ \emph{be cautious about
        reading too much into}; ``on has the same properties'' $\to$ \emph{or has
        the same properties}; ``In all three case'' $\to$ \emph{cases}; ``anwers''
        $\to$ \emph{answers}. The line ``general equilibrium, macro econ, and
        empirical issues'' has no verb; read as the areas where aggregation
        matters.
  \item Page 3: ``The answers to the first question is'' $\to$ \emph{The answer
        \dots\ is}; ``marshallian'' $\to$ \emph{Marshallian} (MWG's name,
        \emph{Walrasian}, added); ``a vector on bundles'' $\to$ \emph{a vector,
        i.e.\ a bundle}; ``With $I$ agent'' $\to$ \emph{agents};
        ``infomation'' $\to$ \emph{information}. Summation limits $i=1,\dots,I$
        added where the original omits them.
  \item Page 4: ``affirmative an SWEn'' $\to$ \emph{an affirmative answer};
        ``resdistribution'' $\to$ \emph{redistribution}; ``wealth effects and the
        same'' $\to$ \emph{are the same}; ``constant across wealth levels''
        kept. The total derivative is written for each good $k$ (the original
        uses the vector $x_i$).
  \item Page 5: ``Sound restrictive'' $\to$ \emph{Sounds restrictive}; ``All
        agent have'' $\to$ \emph{All agents have}; ``quasilinear in the same
        good'' $\to$ \emph{with respect to the same good}; ``two good'' $\to$
        \emph{two goods}; ``radio'' $\to$ \emph{ratio}; ``Homothetic''
        figure label and ``wealth effects are constant'' kept.
  \item Page 6: ``gormen'' $\to$ \emph{Gorman}; ``event agent'' $\to$
        \emph{every agent}; ``how alike agent are'' $\to$ \emph{agents};
        ``we ask an a hand question, we vantel it to work for wealth distri.''
        $\to$ \emph{we asked a hard question: we wanted it to work for every
        wealth distribution}; ``what we must known'' $\to$ \emph{might know};
        ``if it knows for a nice subset'' $\to$ \emph{works for};
        ``Distrinction'' $\to$ \emph{distribution}.
  \item Page 7, first line: ``E though govenment redistrisbate on ownership. of
        agents in this agent'' is partly illegible; read as \emph{e.g.\ through
        government redistribution, or the agents' ownership of assets in this
        economy}. ``satisfies WARP'' $\to$ \emph{satisfy}; ``a minimal
        alsumption on sonsistency property'' $\to$ \emph{a minimal consistency
        assumption}; ``homogenity'' $\to$ \emph{homogeneity}; ``walra's law''
        $\to$ \emph{Walras' law}; ``individual demands all these properties''
        $\to$ \emph{have all these properties}. Homogeneity of degree zero of
        aggregate demand uses that the rule $w_i(p,w)=\alpha_i w$ is homogeneous
        of degree one in $(p,w)$.
  \item Page 8: ``scearios'' $\to$ \emph{scenario}; ``counter examples WARP
        MAY fail'' $\to$ \emph{counterexamples: WARP may fail};
        ``$\alpha_1=\alpha_2=\bar{2}$'' read as $\tfrac12$ (the budget sets are
        labelled $w/2$); ``agent have the same budget set'' $\to$ \emph{both
        agents have}; ``AVR Demand'', ``AVG Demand'' $\to$ \emph{average
        demand}; ``r.p.'' $\to$ \emph{revealed preferred}; ``garenteed'' $\to$
        \emph{guaranteed}; ``the answe affirmative'' $\to$ \emph{the answer
        affirmative}. The red ``WARP is satisfied'' is read as referring to each
        agent individually.
  \item Page 9: ``On the often hand'' $\to$ \emph{On the other hand}.
\end{transcriptnotes}

\end{document}
